\section{A positive proportion of actual zeta zeros have a definite normal velocity}
\label{sec:analytic-arithmetic-transverse-proportion}

All zero counts in this section use positive ordinates and
count multiplicity unless a simple-zero condition is written
explicitly. Put
\[
 N(T)=\sum_{0<\gamma\le T}1,\qquad
 N_s(T)=\#\{\rho=1/2+i\gamma:0<\gamma\le T,
                                      \ \zeta'(\rho)\ne0\}.
\]
Thus a zero of multiplicity $m$ contributes $m$ to $N(T)$ and
nothing to $N_s(T)$ when $m>1$.

\begin{theorem}[NN-A036: a quantitative arithmetic transversality proportion]
\label{thm:analytic-arithmetic-transverse-proportion}
Assume A. At each simple zero, let
\[
 v_\rho=\frac{\rho\zeta(\rho+1)}{\zeta'(\rho)},\qquad
 r_\rho=\frac{|\Re v_\rho|}{|v_\rho|}\in[0,1],\qquad
 c=\frac1{2\zeta(3)}-\frac13>\frac1{15}.
\]
Then
\begin{equation}\label{eq:analytic-normal-mean}
 \liminf_{T\to\infty}\frac1{N(T)}
                 \sum_{\substack{0<\gamma\le T\\\rho\ {
m simple}}}
                           r_\rho\ \ge c.
\end{equation}
In particular, at least a lower proportion $c$ of all zero
occurrences are simple with $\Re v_\rho\ne0$.

For $0<\eta<3c/2$, let
\[
 G_\eta(T)=\#\{\rho\ {
m simple}:0<\gamma\le T,
                              \ |\Re v_\rho|\ge\eta|v_\rho|\}.
\]
The bounds with their two distinct denominators are
\begin{align}
 \liminf_{T\to\infty}\frac{G_\eta(T)}{N(T)}
       &\ge\frac{c-(2/3)\eta}{1-\eta},
         \label{eq:analytic-total-transverse-density}\\
 \liminf_{T\to\infty}\frac{G_\eta(T)}{N_s(T)}
       &\ge\frac{(3/2)c-\eta}{1-\eta}.
         \label{eq:analytic-simple-transverse-density}
\end{align}
For example, taking $\eta=1/30$ gives the convenient rational
consequences
\[
 \liminf\frac{G_{1/30}(T)}{N(T)}\ge\frac4{87},\qquad
 \liminf\frac{G_{1/30}(T)}{N_s(T)}\ge\frac2{29}.
\]
These are lower bounds for the magnitude of the relative
normal component. They assign neither its sign nor a half
proportion to left and right labels, and do not establish
transversality at every simple zero.
\end{theorem}

\begin{lemma}[The exterior Euler phase averaged at the zeros]
\label{lem:analytic-exterior-phase-average}
Under RH, for each fixed $\sigma>1$,
\begin{equation}\label{eq:analytic-exterior-ratio-mean}
 \lim_{T\to\infty}\frac1{N(T)}
       \sum_{0<\gamma\le T}
       \frac{\zeta(\sigma+i\gamma)}
            {\zeta(\sigma-i\gamma)}
                     =\frac1{\zeta(2\sigma)}.
\end{equation}
This is an unweighted zero average with multiplicity, not a
joint independence assertion involving the Hardy phase.
\end{lemma}

\begin{proof}
For every fixed $x>1$, Landau's formula is
\[
 \sum_{0<\gamma\le T}x^\rho
      =-\frac{T}{2\pi}\Lambda(x)+O_x(\log T),
\]
where $\Lambda(x)=\log p$ when $x=p^k$ with integer $k\ge1$,
and is zero for other real $x>1$. A primary locator is
S.~M.~Gonek, \emph{A Formula of Landau and Mean Values of
$\zeta(s)$}, author-hosted PDF, p.~1, equation~(1);
its uniform extension is Theorem~1, equation~(3).
The fixed-$x$ formula is all that is used here.
Under RH, $x^\rho=\sqrt x\,e^{i\gamma\log x}$ and
$N(T)\sim T\log T/(2\pi)$. Thus
\begin{equation}\label{eq:analytic-fixed-frequency-zero-average}
 \frac1{N(T)}\sum_{0<\gamma\le T}e^{i\gamma\log x}\longrightarrow0
                    \qquad(x>0,\ x\ne1).
\end{equation}
For $0<x<1$ use complex conjugation and $1/x>1$.

For $\sigma>1$, the following double series is absolutely and
uniformly convergent in real $t$:
\[
 \frac{\zeta(\sigma+it)}{\zeta(\sigma-it)}
 =\sum_{n,m\ge1}\frac{\mu(m)}{(nm)^\sigma}
                                  e^{it\log(m/n)}.
\]
Indeed the absolute coefficient sum is at most
$\zeta(\sigma)^2<\infty$. Each normalized zero average of a
term has modulus at most one, so dominated convergence for
this summable family applies to
\eqref{eq:analytic-fixed-frequency-zero-average}.
Only $m=n$ survives, and its sum is
$\sum_{n\ge1}\mu(n)n^{-2\sigma}=1/\zeta(2\sigma)$.
This proves \eqref{eq:analytic-exterior-ratio-mean} without
requiring uniformity of Landau's error in $x$.
\end{proof}

\begin{proof}[Proof of Theorem~\ref{thm:analytic-arithmetic-transverse-proportion}]
Let $\theta(t)$ be the usual Hardy--Riemann--Siegel phase and
use the Euler-product branch
\[
 \phi(t)=\arg\zeta(3/2+it),\qquad
 \delta(t)=\arctan\frac1{2t},\qquad
 \psi(t)=\phi(t)-\delta(t).
\]
The first necessary phase average is already a proved
classical statement, not full Gram equidistribution:
\begin{equation}\label{eq:analytic-theta-first-mode}
 \frac1{N(T)}\sum_{0<\gamma\le T}e^{2i\theta(\gamma)}
                                      \longrightarrow0.
\end{equation}
The exact locator is J.~Arias de Reyna,
\emph{On the distribution (mod 1) of the normalized zeros of
the Riemann Zeta-function}, arXiv:1403.7126v1,
Corollary~3.2, equation~(13), with $\kappa=1$; its proof gives
$O(1/\log T)$ in this case. That result sums
$\rho=1/2+i\alpha$ without assuming RH. Here A makes
$\alpha=\gamma$ real, with precisely the multiplicity
convention used in $N(T)$.

Since $e^{2i\phi(t)}=\zeta(3/2+it)/\zeta(3/2-it)$,
Lemma~\ref{lem:analytic-exterior-phase-average} gives
\[
 \frac1{N(T)}\sum_{0<\gamma\le T}e^{2i\phi(\gamma)}
                                  \longrightarrow\frac1{\zeta(3)}.
\]
The replacement of $\phi$ by $\psi$ has vanishing average
error. To see this with no interchange of divergent sums,
fix $Y>0$: the finitely many zeros below $Y$ contribute
$o_T(1)$, while for $\gamma>Y$ the difference of the two
unit phases is at most $2\delta(Y)$. Let $T\to\infty$ and
then $Y\to\infty$. Taking real parts, the two limits are
\begin{equation}\label{eq:analytic-two-separate-real-means}
 \langle\cos2\theta\rangle_T\longrightarrow0,
 \qquad
 \langle\cos2\psi\rangle_T\longrightarrow\frac1{\zeta(3)}.
\end{equation}
No joint distribution or decorrelation has been used.

Define at every zero occurrence, including multiple zeros,
the bounded phase observable
$f(\gamma)=|\cos(\theta(\gamma)+\psi(\gamma))|$.
The pointwise trigonometric identity gives
\[
 \cos2\theta+\cos2\psi
   =2\cos(\theta+\psi)\cos(\theta-\psi)
   \le2|\cos(\theta+\psi)|=2f.
\]
Consequently
\begin{equation}\label{eq:analytic-all-phase-lower}
 \liminf_{T\to\infty}\langle f\rangle_T
                              \ge\frac1{2\zeta(3)}.
\end{equation}
This deterministic inequality is why the separate first
moments suffice for the present lower bound, although they
do not prove the parity-weighted half-label law.

At a simple zero, $Z(t)=e^{i\theta(t)}\zeta(1/2+it)$ is
real, $Z'(\gamma)\ne0$, and
\[
 \zeta'(\rho)=-i e^{-i\theta(\gamma)}Z'(\gamma),\qquad
 v_\rho=\frac{i\rho e^{i\theta(\gamma)}
                   \zeta(3/2+i\gamma)}{Z'(\gamma)}.
\]
For $\gamma>0$, $\arg(i\rho)=\pi-\delta(\gamma)$.
The real sign of $Z'(\gamma)$ only adds a multiple of $\pi$
to this argument. Therefore the corrected exact velocity
formula of NN-A001 gives
\begin{equation}\label{eq:analytic-phase-is-relative-normal}
                   r_\rho=f(\gamma)
                       \qquad(\rho\text{ simple}).
\end{equation}
The Euler product makes the numerator of $v_\rho$ nonzero,
so the displayed normalization is well defined.

The remaining imported input is precisely Montgomery's
proved simple-proportion bound under RH:
\begin{equation}\label{eq:analytic-montgomery-simple-proportion}
                   N_s(T)\ge(2/3+o(1))N(T).
\end{equation}
The source is \emph{The Pair Correlation of Zeros of the
Zeta Function}, Corollary~2, printed p.~182, equation~(7).
Its proof in section~4, pp.~191--192, explicitly states
that the zero sums count multiplicity; a multiplicity-$m$
zero is entered $m$ times and contributes $m^2$ to the
equal-ordinate pair sum. Thus
\eqref{eq:analytic-montgomery-simple-proportion} has exactly
the denominator $N(T)$ used in both phase averages. This is
the already catalogued classical input NT-Z02, not a new
simplicity claim or a full pair-correlation conjecture.

Because $0\le f\le1$, removing the nonsimple occurrences
from \eqref{eq:analytic-all-phase-lower} loses at most
$(N(T)-N_s(T))/N(T)\le1/3+o(1)$. Together with
\eqref{eq:analytic-phase-is-relative-normal}, this proves
\eqref{eq:analytic-normal-mean}. Since $r_\rho\le1$ and
vanishes at a tangent simple zero, it also proves the lower
proportion $c$ of simple transverse zeros.

For the sharper threshold bounds, put
$p_T=N_s(T)/N(T)$ and $g_T=G_\eta(T)/N(T)$. The pointwise
upper bounds on the three subsets (nonsimple, simple below
threshold, and simple above threshold) give
\[
 \langle f\rangle_T
       \le (1-p_T)+\eta(p_T-g_T)+g_T
       =1-(1-\eta)p_T+(1-\eta)g_T.
\]
Hence
\[
 g_T\ge p_T-\frac{1-\langle f\rangle_T}{1-\eta}.
\]
Taking lower limits using $p_T\ge2/3+o(1)$ proves
\eqref{eq:analytic-total-transverse-density}. Dividing the
same inequality by $p_T$, which is bounded away from zero,
and using
\[
 \frac{1-\langle f\rangle_T}{p_T}
       \le\frac{1-1/(2\zeta(3))+o(1)}{2/3+o(1)}
\]
proves \eqref{eq:analytic-simple-transverse-density}.
There is no hidden conversion from a distinct-zero count.

Finally,
\[
 \zeta(3)=1+\frac18+\sum_{n\ge3}\frac1{n^3}
       <1+\frac18+\int_2^\infty x^{-3}\,dx=\frac54,
\]
so $c>2/5-1/3=1/15$. Inserting $\eta=1/30$ in the two
exact bounds gives respectively a number greater than
$4/87$ and a number greater than $2/29$, proving the stated
rational consequences without a numerical constant
calculation.
\end{proof}

\paragraph{Scope.}
The theorem supplies a quantitative lower density of actual
arithmetic simple zeros at which the normal velocity is a
fixed fraction of the full velocity. It uses A, the single
proved Hardy-phase mode, the fixed-frequency Landau formula,
and the classical $2/3$ simplicity proportion. It does not
use a normalized CLT, a joint phase law, full GUE, or an
assumption that all zeros are simple. The positive and
negative signs may remain correlated with index parity;
the half-label law and the exclusion of every simple tangent
remain open here.

\paragraph{AG and harmonic-analysis translation.}
The zero divisor is counted with its actual local lengths.
At a reduced point, the equation $\zeta_t(s)=0$ gives
$d\zeta_t=\zeta'(\rho)\,ds-\rho\zeta(\rho+1)\,dt$;
its section has tangent $\partial_t+v_\rho\partial_s$.
The observable $r_\rho$ is the absolute normalized real
normal component of this tangent relative to the critical
line. The theorem gives a positive lower counting density
on that reduced locus, retaining all multiplicities in the
ambient denominator. The K\"ahler-differential local model
is the one of SGA~1, Expos\'e~I, section~3, Proposition~1.
For the exterior phase, the characters $x\mapsto
e^{i\gamma\log x}$ are characters of the multiplicative
group $\mathbb Q_{>0}$, whose compact dual is the product
of prime circles. Landau's formula identifies the relevant
Haar means, and absolute Dirichlet convergence makes the
ratio's constant term $\sum\mu(n)n^{-3}=1/\zeta(3)$.
This is an actual harmonic-analysis computation, not an
asserted Frobenius equidistribution theorem for the zeros
or an arithmetic family realizing the denominator shift.
